\chapter{Reciprocal-Deficit Normal Form for the Solid PF3 Margin}
\label{chap:g023-reciprocal-deficit}

\status{Exact reparametrization and local sufficient condition; the global reciprocal-deficit law for the actual Riemann coefficient sequence remains open.}

\section{Purpose}

Chapter~\ref{chap:g022-pf3-curvature-barrier} isolates two local conditions
that are sufficient for positivity of the G006 solid order-three minor. This
chapter reparametrizes those conditions by the reciprocal PF$_2$ curvature
deficit. The resulting variables turn the nonlinear G022 barrier into a
one-step Lipschitz bound and put the G006 margin itself into a shifted
multiplicative normal form.

No global PF$_3$ conclusion is drawn here.

\section{Reciprocal curvature deficit}

For the coefficient ratios
\[
 r_k=\frac{a_k}{a_{k-1}},
 \qquad
 q_k=\frac{r_{k+1}}{r_k},
\]
define, whenever $0<q_k<1$,
\[
 \boxed{
 E_k:=\frac{1}{1-q_k}>1.
 }
\]
Thus
\[
 q_k=1-\frac1{E_k}.
\]
At the G006 solid minor indexed by $k$, write
\[
 u=q_{k-1},\qquad v=q_k,\qquad w=q_{k+1}.
\]

\section{Exact normal form}

\begin{theorem}[SOH-G023 reciprocal-deficit normal form]
The G006 ratio-curvature margin
\[
 M_k=(1-v)^2-v^2(1-u)(1-w)
\]
satisfies
\[
 \boxed{
 M_k=
 \frac{E_{k-1}E_{k+1}-(E_k-1)^2}
      {E_{k-1}E_k^2E_{k+1}}.
 }
\]
Consequently the sign of the solid order-three minor is the sign of
\[
 \boxed{
 \widehat M_k:=E_{k-1}E_{k+1}-(E_k-1)^2.
 }
\]
\end{theorem}

\begin{proof}
Use
\[
 1-u=E_{k-1}^{-1},
 \qquad
 1-v=E_k^{-1},
 \qquad
 1-w=E_{k+1}^{-1},
\]
and
\[
 v=\frac{E_k-1}{E_k}.
\]
Substitution into the G006 formula gives the displayed quotient directly. Its
denominator is positive.
\end{proof}

Thus the solid-minor condition becomes
\[
 \boxed{
 E_{k-1}E_{k+1}\ge(E_k-1)^2.
 }
\]

\section{Increment decomposition}

Define
\[
 \alpha_k:=E_k-E_{k-1},
 \qquad
 \beta_k:=E_{k+1}-E_k.
\]
Then
\[
 E_{k-1}=E_k-\alpha_k,
 \qquad
 E_{k+1}=E_k+\beta_k.
\]

\begin{theorem}[SOH-G023 increment decomposition]
The transformed margin satisfies
\[
 \boxed{
 \widehat M_k
 =(E_k-1)+E_k(1-\alpha_k)+E_{k-1}\beta_k.
 }
\]
\end{theorem}

\begin{proof}
Expand
\[
 (E_k-\alpha_k)(E_k+\beta_k)-(E_k-1)^2
\]
and collect terms.
\end{proof}

\begin{corollary}[Local monotone one-Lipschitz certificate]
If
\[
 E_k>1,
 \qquad
 \alpha_k\le1,
 \qquad
 \beta_k\ge0,
\]
then
\[
 \boxed{\widehat M_k>0},
 \qquad
 \boxed{M_k>0}.
\]
\end{corollary}

\begin{proof}
The first term $E_k-1$ is strictly positive. The second is non-negative when
$\alpha_k\le1$, and the third is non-negative when $\beta_k\ge0$ because
$E_{k-1}>0$.
\end{proof}

\section{Equivalence with the G022 local conditions}

The G022 forward-order condition
\[
 q_{k+1}\ge q_k
\]
is equivalent to
\[
 E_{k+1}\ge E_k,
\]
hence to
\[
 \boxed{\beta_k\ge0.}
\]

The G022 barrier
\[
 q_k(2-q_{k-1})\le1
\]
is equivalent to
\[
 \frac1{1-q_k}-\frac1{1-q_{k-1}}\le1,
\]
hence to
\[
 \boxed{\alpha_k\le1.}
\]

Therefore the global law
\[
 \boxed{
 0\le E_{k+1}-E_k\le1
 }
\]
for every relevant index would imply strict positivity of every G006 solid
order-three minor. SOH-G023 does not prove this global law.

\section{Continuous normalized-moment target}

Let
\[
 L(p)=\log R(p),
 \qquad
 R(p)=\Gamma(p+1)^{-1}\int_0^\infty y^p\Phi(y)\,dy,
\]
and define
\[
 \delta(p)=L(p+2)-2L(p)+L(p-2).
\]
At even arguments,
\[
 q_k=e^{\delta(2k)}.
\]
Introduce
\[
 \boxed{
 \mathcal E(p)=\frac{1}{1-e^{\delta(p)}}.
 }
\]
Then $E_k=\mathcal E(2k)$. A sufficient continuous route to the global G023
increment law is
\[
 \boxed{
 0\le\mathcal E'(p)\le\frac12.
 }
\]
Indeed, integrating over an interval of length two gives
\[
 0\le E_{k+1}-E_k\le1.
\]

Differentiation yields
\[
 \mathcal E'(p)
 =\frac{e^{\delta(p)}\delta'(p)}{(1-e^{\delta(p)})^2}.
\]
Hence the continuous derivative package is
\[
 \boxed{\delta'(p)\ge0}
\]
together with
\[
 \boxed{
 \delta'(p)\le\cosh(\delta(p))-1.
 }
\]
These inequalities are a new explicit analytic target, not a proved theorem in
this chapter.

\section{Finite diagnostic}

The accompanying receipt computes $E_k$ and consecutive increments directly
from finite high-precision positive-kernel moments. On the sampled range the
increments lie strictly between zero and one and decrease numerically. This is
recorded only as \texttt{FINITE\_DIAGNOSTIC\_NOT\_PROOF}.

\section{Compatibility with G021}

For the G021 counterexample
\[
 u=\frac15,
 \qquad
 v=\frac9{10},
 \qquad
 w=\frac15,
\]
one obtains
\[
 E_{k-1}=\frac54,
 \qquad
 E_k=10,
 \qquad
 E_{k+1}=\frac54.
\]
Therefore
\[
 \alpha_k=\frac{35}{4}>1,
 \qquad
 \beta_k=-\frac{35}{4}<0.
\]
The counterexample lies outside the G023 sufficient package in both directions.

\section{Open obligations}

The active reciprocal-deficit proof targets are now:
\begin{enumerate}
 \item prove $E_{k+1}\ge E_k$ for every actual Riemann coefficient index;
 \item prove $E_{k+1}-E_k\le1$ for every such index;
 \item equivalently, prove a sufficient continuous bound such as
       $0\le\mathcal E'(p)\le1/2$;
 \item after all solid minors are controlled, provide the separate upgrade to
       all order-three Toeplitz minors required for PF$_3$.
\end{enumerate}

\gap{The global monotone one-Lipschitz reciprocal-deficit law is not proved. Therefore all solid minors globally, full PF$_3$, PF$_\infty$, SOH-G003, SOH-C005, and RH remain open.}
